% ===========================================================================
%  LOOM -- Overleaf edition · demo / starter document
%  Compile with XeLaTeX (Menu > Compiler > XeLaTeX on Overleaf).
%  Nothing to upload: every font ships with Overleaf's TeX Live.
%
%  The [cjk] option turns on Chinese via Noto CJK SC (bundled on Overleaf;
%  macOS Songti/Kaiti as a local fallback).  Drop it for a pure-English doc.
% ===========================================================================
\documentclass[cjk]{loom}

\runningthread{loom · overleaf edition}

\begin{document}

\loomcover
  {Loom}
  {A woven notebook class · Overleaf edition}
  {北极甜虾 (剑心犹在！)}
  {\today}

{\small\tableofcontents}
\clearpage

% ===========================================================================
\section{The weave at a glance}
\warmth{3}

\begin{strand}
A page is cloth on a loom: you don't just typeset the maths, you typeset the
\keyword{weave} of your own understanding. The \emph{warp} runs along the left
selvage rail, a \emph{knot} binds threads into a result, and a \emph{loose}
thread is a question you haven't closed yet.
\end{strand}

This edition keeps every command of the local \textsc{loom} class but swaps the
macOS display faces for \texttt{Libertinus Sans}, so it compiles on Overleaf
with \keyword{no font uploads}. The body is Libertinus Serif and the maths is
Libertinus Math, exactly as in the original.\loose{Try editing the dye colours
in \texttt{loom.cls} -- indigo, madder, weld.}

\begin{strand}
中英混排也开箱即用：用 \texttt{[cjk]} 选项即可。中文走 Noto CJK SC（Overleaf
自带），\emph{强调}的中文在本地回退到楷体。一行公式 $\int_0^1 x^2\,dx = \tfrac13$
照常排版。
\end{strand}

\subsection{Knots} A knot is a result that binds threads together.

\begin{definition}[Woven space]
A \keyword{woven space} is a set $W$ together with a family of threads
$\{t_i\}_{i\in I}$, $t_i \subseteq W$, whose union is all of $W$.
\end{definition}

\begin{theorem}[Selvage bound]
For any finite woven space $W$ with $n$ threads, the number of knots is at most
$\binom{n}{2}$, with equality iff every pair of threads crosses exactly once.
\end{theorem}

\begin{proof}
Each knot is determined by the unordered pair of threads it binds, so the count
is bounded by the number of such pairs, $\binom{n}{2}$. Equality forces every
pair to meet, which is the stated crossing condition.
\end{proof}

\begin{example}
The complete weave on $4$ threads has $\binom{4}{2}=6$ knots.
\end{example}

% ===========================================================================
\section{Fill-in study notes}
\warmth{1}

The pedagogy tools turn a passive read into active recall: state the result, but
leave the load-bearing steps \keyword{blank}.

\begin{definition}[Continuity, fill the clause]
$f\colon X\to Y$ is \keyword{continuous} at $x_0$ if for every
$\varepsilon>0$ there is a $\delta>0$ such that
\[ d_X(x,x_0)<\delta \implies d_Y\big(f(x),f(x_0)\big) < \fillin[1.4cm]. \]
\end{definition}

\begin{proof}[Sketch, with a gap]
Compose the two estimates and chase the quantifiers.
\TODO{which $\delta$ makes the implication hold?}
\end{proof}

\recall{State the $\varepsilon$--$\delta$ clause without looking.}

\begin{yourturn}
Prove that the composition of two continuous maps is continuous. Use the
workspace below.
\end{yourturn}
\workspace[3]

% ===========================================================================
\section{Margins, warps and warmth}
\warmth{5}

A \loose{is this tight enough?} loose thread parks an open question in the
margin. A \warp{key lemma} declares a recurring thread you can \pick{key lemma}
later. The warmth gauge \warmth{4} records how well you actually grok a block,
from \warmth{0} to \warmth{5}.

\begin{remark}
Everything here is plain XeLaTeX. Delete the demo content, keep
\texttt{loom.cls}, and start weaving your own notes.
\end{remark}

\end{document}
